\documentclass[10pt]{beamer}
\usetheme[progressbar=frametitle]{metropolis}
\usepackage{tikz}
\usetikzlibrary{positioning,arrows,shapes,fit,calc}
\usepackage{appendixnumberbeamer}
\usepackage[svgnames,table]{xcolor}
\usepackage{tabularx}
\usepackage{booktabs}
\usepackage[scale=2]{ccicons}
\usepackage{adjustbox}
\usepackage{pifont}
\usepackage{pgfplots}
\usepgfplotslibrary{dateplot}
\usepackage{xspace}
\usepackage{amsmath,amssymb}

\newcommand{\themename}{\textsc{metropolis}\xspace}

% Math commands
\newcommand{\E}{\mathbb{E}}
\renewcommand{\max}{\text{max}}
\renewcommand{\min}{\text{min}}
\newcommand{\argmax}{\text{argmax}}

\begin{document}

%==============================================================
% Stochastic Games - More Detail
%==============================================================

\begin{frame}{Multi-Agent Settings: Stochastic Games (Detail)}
    \begin{itemize}
        \item Stochastic game: $(S, A_1, \ldots, A_n, P, R_1, \ldots, R_n, \gamma)$
        \begin{itemize}
            \item Each stage is a matrix game determined by state $s$
            \item Transitions depend on joint actions of all players
        \end{itemize}
        \vspace{0.2cm}
        \item Value function for player $i$ given joint policy $\pi = (\pi_1, \ldots, \pi_n)$:
        \begin{equation}
            V_i^{\pi}(s) = \E_{\pi}\left[\sum_{t=0}^{\infty} \gamma^t R_i(s_t, a_t^1, \ldots, a_t^n) \,\middle|\, s_0 = s\right]
        \end{equation}
        \vspace{0.2cm}
        \item Nash Equilibrium in stochastic game:
        \begin{equation}
            V_i^{\pi^*}(s) \geq V_i^{(\pi_i, \pi_{-i}^*)}(s) \quad \forall i, \forall s, \forall \pi_i
        \end{equation}
    \end{itemize}
\end{frame}

%==============================================================
% Nash Q-Learning
%==============================================================

\begin{frame}{Multi-Agent RL: Nash Q-Learning}
    \begin{itemize}
        \item Extension of Q-learning to stochastic games:
        \begin{equation}
            \hat{Q}_i(s,\mathbf{a}) \leftarrow \hat{Q}_i(s,\mathbf{a}) + \alpha \left[ r_i + \gamma \cdot NashQ_i(s',\hat{Q}_1,\ldots,\hat{Q}_n) - \hat{Q}_i(s,\mathbf{a}) \right]
        \end{equation}
        where $\mathbf{a} = (a_1,\ldots,a_n)$ and $NashQ_i(s,\hat{Q}_1,\ldots,\hat{Q}_n)$ is player $i$'s expected payoff in state $s$ under the Nash equilibrium
        \vspace{0.2cm}
        \item Challenges:
        \begin{itemize}
            \item Computing Nash equilibrium in each state
            \item Multiple equilibria problem
            \item Non-stationarity during learning
        \end{itemize}
        \vspace{0.2cm}
        \item Convergence guarantees:
        \begin{itemize}
            \item Requires unique Nash equilibrium at each state
            \item Convergence to Nash Q-values under appropriate conditions
        \end{itemize}
    \end{itemize}
\end{frame}

%==============================================================
% Perfect Bayesian Equilibria
%==============================================================

\begin{frame}{Multi-Agent Settings: Perfect Bayesian Equilibria}
    \begin{itemize}
        \item Bayesian game with private information:
        \begin{itemize}
            \item Players have types $\theta_i \in \Theta_i$ with prior distribution $\mu(\theta)$
            \item Strategy maps types to actions: $\sigma_i: \Theta_i \rightarrow \Delta(A_i)$
        \end{itemize}
        \vspace{0.2cm}
        \item Perfect Bayesian Equilibrium (PBE) requires:
        \begin{enumerate}
            \item Sequential rationality: Each player's strategy maximizes expected utility given beliefs at every information set
            \item Consistent beliefs: Updated via Bayes rule whenever possible
        \end{enumerate}
        \vspace{0.2cm}
        \item In sequential settings:
        \begin{itemize}
            \item Information revealed over time changes beliefs
            \item Strategies at later stages depend on updated beliefs
            \item Actions may signal private information
        \end{itemize}
    \end{itemize}
\end{frame}

%==============================================================
% Neural Fictitious Self-Play
%==============================================================

\begin{frame}{Learning in Bayesian Games: Neural Fictitious Self-Play}
    \begin{itemize}
        \item NFSP combines:
        \begin{enumerate}
            \item RL Component: Learn best response to opponents
            \item SL Component: Track average strategy via supervised learning
        \end{enumerate}
        \vspace{0.2cm}
        \item Key components:
        \begin{itemize}
            \item Best response network: $\hat{Q}_i(s,a;\theta_i)$ updated via DQN
            \begin{equation}
                L(\theta_i) = \E_{(s,a,r,s')} \left[ \left( r + \gamma \max_{a'} \hat{Q}_i(s',a';\theta_i^-) - \hat{Q}_i(s,a;\theta_i) \right)^2 \right]
            \end{equation}
            
            \item Average strategy network: $\bar{\pi}_i(a|s;\phi_i)$ trained via supervised learning
            \begin{equation}
                L(\phi_i) = -\E_{(s,a)} \left[ \log \bar{\pi}_i(a|s;\phi_i) \right]
            \end{equation}
        \end{itemize}
        \vspace{0.2cm}
        \item Converges to approximate Nash equilibrium in various games
    \end{itemize}
\end{frame}

%==============================================================
% Korean Auction - More Detail
%==============================================================

\begin{frame}{Korean Auction: Perfect Bayesian Equilibrium}
    \begin{itemize}
        \item Strategic considerations:
        \begin{itemize}
            \item First round: Signal value vs. conceal information
            \item Second round: Update beliefs based on first-round outcome
        \end{itemize}
        \vspace{0.2cm}
        \item Belief updating:
        \begin{itemize}
            \item Signal $x_i = 1$: Opponent's value likely below threshold
            \item Signal $x_i = 0$: Opponent's value likely above threshold
        \end{itemize}
        \vspace{0.2cm}
        \item Equilibrium bidding:
        \begin{itemize}
            \item First round: Shading below true value
            \item Second round (if won first round): Bid lower based on updated belief
            \item Second round (if lost first round): Bid higher if value sufficiently high
        \end{itemize}
        \vspace{0.2cm}
        \item RL approach:
        \begin{itemize}
            \item Separate neural networks for each round
            \item Condition second-round network on first-round outcome
        \end{itemize}
    \end{itemize}
\end{frame}

%==============================================================
% Entry-Exit Game - More Detail
%==============================================================

\begin{frame}{Markov Perfect Equilibrium in Entry-Exit Games}
    \begin{itemize}
        \item Model based on Ericson \& Pakes (1995):
        \begin{itemize}
            \item Firms decide to enter, stay, or exit based on market state
            \item Market evolves based on exogenous shocks and firm decisions
            \item Profits depend on market structure (competition effects)
        \end{itemize}
        \vspace{0.2cm}
        \item Markov Perfect Equilibrium computation:
        \begin{enumerate}
            \item For each state $s$, compute stage game using continuation values
            \item Find Nash equilibrium of the stage game
            \item Update value functions according to equilibrium strategies
            \item Iterate until convergence
        \end{enumerate}
        \vspace{0.2cm}
        \item Computational challenges:
        \begin{itemize}
            \item Curse of dimensionality with many firms or state variables
            \item Multiple equilibria (equilibrium selection)
            \item Numerical stability issues
        \end{itemize}
    \end{itemize}
\end{frame}

\end{document}